\section{Introduction}
\label{sec:intro}

%talk about representation learning and why it is important
Representation Learning is a fundamental problem in the field of deep learning. Most modern deep learning models follow a two-step process: 1) Pre-training: train models to extract powerful representations from raw data, and 2) Fine-tuning: use the representations to learn models that solve a specific downstream task. The pre-training stage is also known as Representation Learning, as the fundamental goal of this stage is to learn a transformation that converts raw input into meaningful representations in some vector space. Many state-of-the-art models in various domains, including image classification~\cite{oquab2024dinov2}, text-to-image generation~\cite{ramesh2022hierarchical}, and image segmentation~\cite{oquab2024dinov2}, have demonstrated that proper pre-training is crucial for achieving good performance on the desired downstream task. As a result, developing effective algorithms for pre-training has become a critical problem in deep learning. 

\subsection{Motivation}
While pre-training has proven highly effective, many strategies are often transferred directly across different data modalities. For example, SimCLR~\cite{chen2020simple} introduced contrastive pretraining for images, which was later adopted for audio~\cite{elizalde2023clap}, video~\cite{sun2022long}, and medical images~\cite{wang2022medclip} with minimal adaptation. Although a direct adaptation approach can yield reasonable results, it frequently overlooks opportunities to incorporate modality-specific properties that could enrich the learned representations. Moreover, pre-training algorithms sometimes have assumptions about the input data that may not hold in a new modality. For example, masked modeling as a pre-training strategy proved to be highly effective in the language domain but remained relatively inferior in the image domain. One key reason for this was the difference in information density between the two data modalities. Language, being a human-generated signal, is highly information dense, and predicting a few missing words required models to develop strong language understanding. Images, on the other hand, are natural signals with huge redundancy, and therefore, predicting a few missing image patches becomes a trivial task that can be solved without the model developing any visual understanding. Thus, a key motivation when designing pre-training methods should be not only to generalize across modalities, but also to exploit the unique structures and inductive biases inherent to each data type.% Why are we motivated to automate this problem, maybe compare with existing methods to the proposed one, maybe also availability of large scale imagery and crowd-sourced geotagged data.

%gap in existing literature?

%The availability of overhead imagery for any location in the world, along with geotagged audios hosted on various open-sourced platforms, enables us to address such a multi-modal correspondence problem and design appropriate methods.


\subsection{Problem Statement}
% Problem formulation, primary challenges to address
We propose developing data-aware representation learning methods for geospatial data. Learning representations for geospatial data can be formulated as a problem of learning a neural field over the earth. More specifically, we want to learn an implicit representation for the entire earth, such that, for any point on earth $(lat,lon)$, we learn a function $R=F_\theta(lat,lon)$, where $R$ is a vector representation of the query location. Depending on the training design, $R$ can capture different information about the query location, such as visual semantics, text descriptions, ecological properties, etc. In this disseratation, first, we will develop a framework to learn general-purpose neural representation field over the earth. We will examine the existing representation learning strategies for geospatial coordinates and demonstrate existing methods are incompatible with coordinate data. To that end, we will propose a new framework that addresses the existing limitations. Second, we will develop a specialized neural field of text representations using satellite images as a proxy for location. We will exploit the cross-view relationship between satellite images and ground-level images, enabling us to train this model without requiring any text data.  Finally, we will examine the potential benefits and risks of geospatial generative models from the perspective of representation learning. We will design methods to better understand the risks associated with current generative models and propose solutions to mitigate those risks. We will utilize the findings to enhance geospatial generative models and leverage them to further refine representation learning pipelines. In summary, my dissertation will broadly be aimed at the following:
\begin{itemize}
     \item Propose a general-purpose neural representation field for geolocation. Understand the limitations of directly adapting existing methods.
    \item Learn a specialized neural field of textual representations using satellite images as input. Use modality-specific properties to circumvent issues. 
    \item Understand the potential risks and benefits of geospatial generative models for representation learning. 
    \begin{itemize}
        \item Propose a new metric to evaluate existing geospatial generative models.
        \item Propose SimLBR, a new fake image detection strategy that learns a tight decision boundary around the real image distribution, enabling detection of generated images from unseen generators in real-world scenarios.
        \item Use the findings of deepfake detection experiments to improve geospatial generative models. Use the improved generative models to generate synthetic data and scale representation learning pipelines.
    \end{itemize}
\end{itemize}

